Traditional software engineering uses the mechanical abstractions of a programming language to realize the meaning of a software system. The techniques presented in this document take a different approach, they specify the meaning of the software system and use it to automatically construct the closed semantic world as a computable semantic model from the meaning. Since the computable semantic model is built from computable semantic primitives, it can be directly synthesized into an implementation that has no semantic gap with the meaning of the system.

An engine is used to reason abut the CSM's meaning to identify if it is missing meaning or is internally inconsistent. It automatically diagnosis failures by identifying meaningful questions that can't be answered through the CSM and motivates root cause analysis to learn new semantics. In addition to internal inconsistencies, the design can also expand its semantics by enabling the engine to automatically diagnose failures in execution of the CSM. This is done by instrumenting the synthesized program to gather necessary diagnostic data, preserving it at scale using domain specific compression and replaying the execution to allow the engine to reason about it. In the process the design learns new semantics and becomes progressively more intelligent as it evolves to continue realizing its intentions in its operational environment.

As is evident, in this framework, meaning is the primary attribute that drives the entire process. Any process that reasons about the meaning and makes decisions is considered to be intelligent and in this case, the narratives established with the semantic toolkit contain the intentions of intelligence that reasons about the world during execution and makes decisions. This makes software systems inherently possess deterministic intelligence because their meaning is defined through a fixed set of narratives that reason about the world to achieve their intentions. Further more, the engine itself is intelligent because it reasons about the worlds meaning to perform its function. 

In this process, the development of software systems is automated by using a semantic toolkit to automatically build the closed semantic world and its implementation from the specified meaning. The management of software systems is automated by reasoning about the world to automatically diagnose failures and ensure its correctness by construction. Every meaningful question about the world has been asked and answered, enabling the design to learn new meaningful questions through root cause analysis on failures. Through this process of exhausting every meaningful question and answering it, the software system is managed by construction.

Finally, as a result of the processes described in this document, the software systems that run modern society are less likely to fail and if they do fail, an automated process ensures their recovery and automated failure diagnosis supports deterministic learning to prevent future failures. This framework transforms the construction of software systems into a process of meaningful world building while using an engine to reason about its correctness. The future of software systems will be less about writing bug free code and more about building intelligent and secure systems that understand each other. 